\section{Quick start}
\label{sec:quickstart}

\subsection{Solve a scramble you were given}

\begin{lstlisting}[style=shell]
uv run rubiks-cfop --scramble "R U R' U' F2 D"
\end{lstlisting}

The program prints the scrambled cube as an unfolded net in real colors, then the
solution stage by stage.  Piped into a file, or with \texttt{--no-color}, the net is
drawn with colored squares instead of terminal colors.  A run looks like this:

\begin{lstlisting}[style=shell]
           Top
        ### ### ###
        ### ### ###    <- U face (yellow when the cross is white)
        ### ### ###
  Left    Front    Right    Back
### ###  ### ###  ### ###  ### ###  ...
### ###  ### ###  ### ###  ### ###  ...   <- L, F, R and B in a band
### ###  ### ###  ### ###  ### ###  ...
           Down
        ### ### ###
        ### ### ###    <- D face (the cross color)
        ### ### ###

Orientation: none
Cross color: yellow (bottom layer)
Scramble: R U R' U' F2 D

Cross: D' F2
F2L1: R' F R F'
F2L2: (none)
F2L3: (none)
F2L4: (none)
OLL (37): F R' F' R U R U' R'
PLL (skip): (none)

Total: 14 moves
\end{lstlisting}

(The \texttt{\#} characters stand for colored cells; Section~\ref{sec:output}
shows the real thing.)

\subsection{Generate a scramble and solve it}

\begin{lstlisting}[style=shell]
uv run rubiks-cfop --random --seed 42
\end{lstlisting}

\texttt{--random} is the default when no scramble is given.  The scramble is a
random-state scramble (WCA 4b3), typically 20 to 22 moves, and \texttt{--seed}
makes it reproducible: the same seed always yields the same scramble.

\begin{lstlisting}[style=shell]
Orientation: zz
Cross color: white (bottom layer)
Scramble: D B2 F2 L2 U' F2 U B2 R2 U L2 F R' U' B' U L R F D B

Cross: z z R' F D' F2 R' F
F2L1: U2 B' R2 B R2
F2L2: F' L F L2 U' L
F2L3: U L U' B U B' L'
F2L4: U B U B' U2 R' U2 R
OLL (40): U R' F R U R' U' F' U R
PLL (F Perm): U' L R2 U R U R2 U' R' U' R2 U' R U2 L' U R' U2

Total: 62 moves
\end{lstlisting}

\subsection{Pick the cross yourself}

\begin{lstlisting}[style=shell]
uv run rubiks-cfop --scramble "R U R' U'" --cross white
uv run rubiks-cfop --scramble "R U R' U'" --cross-options
\end{lstlisting}

\texttt{--cross} names one of the six colors; \texttt{--cross-options} first lists
what each of the six crosses would cost, shortest first:

\begin{lstlisting}[style=shell]
  yellow  0 moves
  white   3 moves
  green   4 moves
  ...
\end{lstlisting}

\subsection{Get the result as JSON}

\begin{lstlisting}[style=shell]
uv run rubiks-cfop --scramble "R U R' U'" --json
\end{lstlisting}

The payload carries the scramble, the cross color, the rotation applied, every
stage with its moves and the facelets \emph{after} the stage, the metrics and a
verification flag.  It is the same structure the Python API returns from
\texttt{Solution.to\_dict()}, so a GUI or a web front-end can consume it directly;
Section~\ref{sec:library} documents it.
